\scanpage{119}
\mspage{2.2}
complete, $\widehat{S^f}\longrightarrow\widehat S^f$ is an isomorphism, if and only if $\widehat{S^f}/\underline m^i\widehat{S^f}\longrightarrow
\widehat S^f/\underline m^i\widehat S^f$ is an isomorphism for all $i$. However, $\widehat{S^f}/\underline m^i\widehat{S^f}=S^f/\underline m^iS^f=
(S/\underline m^iS)^f$ and $\widehat S^f/\underline m^i\widehat S^f=(\widehat S/\underline m^i\widehat S)^f=(S/\underline m^iS)^f$. Consequently $\widehat{S^f}\longrightarrow\widehat S^f$ is
an isomorphism.

Lemma 2.3. Let $S$ be a flat local $R$-algebra and $\underline q$ an ideal of definition
of $R$. Then $e_{S/R}(\underline q)=L_S(S/\underline m S)^{(2)}$ where $\underline m$ is the maximal ideal of $R$ and
where $e_{S/R}(\underline q)=e_S(\underline q S)/e_R(\underline q)$ by definition. If $S$ is also an unramified
$R$-algebra, then $e_S(\underline q S)=e_R(\underline q)$.

\uline{Proof}: Since $S$ is $R$-flat we have that if $R\supset a_1\supset\cdots\supset\underline a_t=q^h$
is a composition series with successive factors isomorphic to $R/\underline m$, then
$S\supset S\otimes\underline a_1\supset\cdots\supset S\otimes a_t=Sq^h$ and the successive factors are isomorphic
to $S\otimes_R\overline R$. Therefore we have that $L_S(S/\underline q^hS)=L_R(R/\underline q^h)L_S(S\otimes_R\overline R)$ which
gives us the result that $e_S(\underline q S)=e_R(q)L_S(S/\underline m S)$ or $e_{S/R}(\underline q)=L_S(S/\underline m S)$. The
rest of the lemma is obvious.

Proposition 2.4. Let $S$ be a local $R$-algebra and $T$ a local $S$-algebra.
Then we have the following:

a) If $S$ is an unramified $R$-algebra, then $r(S/R)=1$ and $r(T/R)=r(T/S)$.

b) If $T$ is an unramified and flat $S$-algebra, then $r(T/R)=r(S/R)$.

c) If $\widehat S$ and $\widehat R$ are the completions of $S$ and $R$ respectively, then
$r(S/R)=r(\widehat S/\widehat R)$.

\uline{Proof}: a) If $S$ is an unramified $R$-algebra, we have that $T_R^f\approx T_S^f$
by lemma 2.1. Therefore $r(T/R)=e(T_R^f/T)=e(T_S^f/T)=r(T/S)$. If $S$ is an
unramified $R$-algebra, then $S_R^f\approx S$ by lemma 2.1, and hence $r(S/R)=e(S_R^f/S)=
e(S/S)=1$.

b) The fact that $T$ is a flat and unramified $S$-algebra entails that
$T_R^e$ is a flat and unramified $S_R^e$-algebra ([6, proposition 1.5]) and hence
\mspage{118}
